\subsection{Medición de $L$}

Se colocó una resistencia $R_s = 470\,\text{m}\Omega$ en serie con el inductor. La corriente
se obtuvo como $i_L = v_{R_s}/R_s$ y la tensión sobre la bobina con el canal diferencial
del osciloscopio ($\text{MATH} = \text{Ch4} - \text{Ch3}$).

En la fase ON la tensión es aproximadamente constante e igual a $V_d$, por lo que
\begin{equation}
    L = \frac{\overline{v_L}|_{\text{ON}}}{di_L/dt},
\end{equation}
donde $di_L/dt$ se obtuvo ajustando una recta a $i_L(t)$ en cada flanco ascendente
($T_{\text{on}} \approx 100{,}5\,\mu\text{s}$, $f = 5\,\text{kHz}$, $D = 0{,}5$).

\begin{table}[h!]
\centering
\caption{Parámetros medidos para la estimación de $L$.}
\label{tab:inductor_medicion}
\begin{tabular}{lcc}
\hline
\textbf{Magnitud} & \textbf{Símbolo} & \textbf{Valor} \\
\hline
Tensión media (fase ON)            & $\overline{v_L}|_{\text{ON}}$ & $3{,}04\,\text{V}$                  \\
Pendiente de corriente             & $di_L/dt$                     & $2{,}96 \times 10^3\,\text{A/s}$   \\
Inductancia medida                 & $L_{\text{med}}$              & $1{,}03\,\text{mH}$                 \\
Inductancia nominal                & $L_{\text{nom}}$              & $1{,}02\,\text{mH}$                 \\
Error relativo                     & $\epsilon$                    & $\approx 1\%$                       \\
\hline
\end{tabular}
\end{table}

El valor obtenido concuerda con el nominal con un error de $\approx 1\%$.